% ---------------------------------------------------------------
% Section: Detection & Orientation Improvements
% ---------------------------------------------------------------

\section[Detection Improvements]{Detection \& Orientation Improvements}

\begin{frame}{Where We Left Off (Exp \#290)}
\begin{columns}
\begin{column}{0.55\textwidth}
\textbf{Status after Exp \#290:}
\begin{itemize}\small
    \item LodeSTAR detects JP positions well
    \item Orientation CNN: works with refined centres
    \item Velocity-vector orientation: 85\% coverage
    \item Single crop used for training
    \item Tested on 100-frame subset only
\end{itemize}
\end{column}
\begin{column}{0.45\textwidth}
\textbf{Open issues:}
\begin{itemize}\small
    \item Low recall on full dataset ($\sim$0.72)
    \item CNN orientation requires labelled data
    \item No cross-frame identity (no tracking)
    \item No quantitative physics output
\end{itemize}
\end{column}
\end{columns}
\end{frame}

\begin{frame}{Multi-Crop Training}
\begin{columns}
\begin{column}{0.55\textwidth}
\textbf{Problem:}
\begin{itemize}\small
    \item Single reference crop → over-fits one appearance
    \item JP particles vary in orientation and lighting
    \item Recall 0.72 — many particles missed
\end{itemize}
\smallskip
\textbf{Solution:}
\begin{itemize}\small
    \item Pool of \textbf{15 cropped particle images} in\\
          \texttt{data/Samples/JP\_Fe\_wf\_2\_40/Samples/}
    \item Each training step samples randomly via \texttt{dt.Value(lambda)}
    \item Self-supervised LodeSTAR objective unchanged
\end{itemize}
\end{column}
\begin{column}{0.45\textwidth}
\begin{center}
\small\textbf{Recall improvement:}\\[0.3cm]
\begin{tabular}{|c|c|}
\hline
\textbf{Crops} & \textbf{Recall} \\
\hline
1 (single) & 0.72 \\
\hline
15 (pool)  & \textbf{0.97} \\
\hline
\end{tabular}
\end{center}
\smallskip
\begin{mynote}[Key insight]
\small Diverse crop pool exposes the model to the full range of appearances without changing the training objective.
\end{mynote}
\end{column}
\end{columns}
\end{frame}

\begin{frame}{Orientation via Template-Matching NCC}
\begin{columns}
\begin{column}{0.50\textwidth}
\textbf{Pipeline:}
\begin{enumerate}\footnotesize
    \item \textbf{Build template bank}:
          load reference crop, apply annular mask,
          pre-rotate at 2° steps (0–360°)
    \item \textbf{Refine position}:
          HoughCircles, fallback to CoM
    \item \textbf{Best-angle match}:
          $\hat{\varphi} = \arg\max_\theta\, \mathrm{NCC}(\text{crop}, T_\theta)$
    \item \textbf{Output}: \texttt{phi}, \texttt{orientation\_ncc} in CSV
\end{enumerate}
\end{column}
\begin{column}{0.50\textwidth}
\textbf{Advantages over CNN:}
\begin{itemize}\footnotesize
    \item No labelled training data
    \item One template crop → instant deployment
    \item NCC score = per-detection confidence
    \item Trivially generalises to new particle types
\end{itemize}
\begin{mynote}
\footnotesize\texttt{--detection-mode template}\\
\texttt{--orientation-template phi245.9.png}
\end{mynote}
\end{column}
\end{columns}
\end{frame}

\begin{frame}{Full Dataset Test Results (JP\_Fe\_wf\_2\_40)}
\begin{columns}
\begin{column}{0.50\textwidth}
\begin{itemize}\small
    \item 11 videos $\times$ 100 frames = 1{,}100 frames total
    \item Widefield microscopy, 30 fps, cutoff = 0.8
    \item Ground truth from YOLO tracker
\end{itemize}
\begin{center}
\small\begin{tabular}{|l|c|}
\hline
\textbf{Metric} & \textbf{Value} \\
\hline
Precision & 0.278 \\
\hline
Recall    & \textbf{0.822} \\
\hline
F1        & 0.457 \\
\hline
\end{tabular}
\end{center}
\end{column}
\begin{column}{0.50\textwidth}
\begin{myalert}[Note on low precision]
\footnotesize GT from YOLO labels only centre-frame particles — many real detections counted as FP. \textbf{Recall is the meaningful metric}: 82\% of labelled particles found.
\end{myalert}
\smallskip
\textbf{Detection modes:}
\begin{itemize}\small
    \item \texttt{constant} — local maxima (baseline)
    \item \texttt{area} — blob area threshold
    \item \texttt{watershed} — touching particles
    \item \texttt{template} — NCC + orientation
\end{itemize}
\end{column}
\end{columns}
\end{frame}
